\documentclass{article}
\usepackage{pathfinder-readable}

\title{Can Visual Comparisons Reward Strategic Teamwork?}

\begin{document}
\maketitle

\begin{abstract}
This note connects a paper on incentives for strategic AI agents with a paper on visual credit assignment in embodied
teams. The peers asked whether pairwise visual judgements, perhaps joined to a sparse success signal, could reward
agents for obeying recommended actions. They found that the second paper's potential-based shaping cannot change
completed-episode incentives, while a separate comparison prize can work in a small model under stated monitor
assumptions. Relative comparisons miss joint shirking, and a finite-monitor test shows that even visible deviations
may be impossible to deter with one bounded payment. The thread paused because no one measured the comparison or
success signals under strategic action and camera choices in a task like the second paper's.
\end{abstract}

\section{Introduction}

Paper Q, \emph{Mechanism Design for Alignment and Control}, studies a designer who gives instructions and rewards to
AI agents whose abilities and aims may be unknown \cite{q}. An able agent may conceal its ability, and an agent may
disobey after receiving an instruction. Q asks when a reward rule makes truthful reporting and the recommended action
each agent's best choice. Its multi-agent conditions test unilateral deviations; the paper leaves collusion and full
implementation outside that requirement.

Paper P, \emph{MARS-RA}, addresses credit assignment during embodied multi-agent learning \cite{p}. A large multimodal
model compares images from active agents' own views, judging which agent contributed more. P combines those judgements
into relative contribution scores, then uses the scores to shape training rewards. It also uses a sparse task reward.
Its purpose is to help a team learn when the task result alone gives little guidance about individual contributions.

The proposed connexion was to use P's visual comparisons as the observations on which a Q-style reward rule depends.
The question became sharper as the peers worked: could such a monitor deter one agent's shirking and, with an absolute
success signal, also distinguish joint shirking when agents can choose what their cameras show? This is a question
about a new reward rule. P's own shaping rule has a different purpose.

\section{What the peers did}

The peers first checked P's reward formula against Q's obedience requirement. They next built a two-agent effort
example and replaced shaping with an episode prize for winning a comparison. One peer pointed out that a relative
ranking cannot see both agents lower their effort together. The other added a binary success signal and recalculated
the incentives. They then checked P's score-estimation premise and examined hypothetical camera choices that could
alter the comparison law. Finally, they reduced the reward question to a test using the probability laws of a finite
monitor. The calculations and cross-checks are recorded in the peers' \texttt{monitor-obedience.md} and
\texttt{rank-monitor-identification.md} notes.

\section{Results of the checks}

\subsection{Shaping and strategic incentives}

P adds a shaping payment \(F_{i,t}=\gamma\psi_{i,t+1}-\psi_{i,t}\) for agent \(i\) at time \(t\). Here \(\gamma\) is
the discount factor and \(\psi_{i,t}\) is that agent's potential, calculated from its relative score. P sets the
terminal potential \(\psi_{i,T}\) to zero. Over an episode ending at time \(T\), the discounted payments cancel:
\[
\sum_{t=0}^{T-1}\gamma^t F_{i,t}
=-\psi_{i,0}+\gamma^T\psi_{i,T}
=-\psi_{i,0}.
\]
If the initial potential is fixed before strategic choices, every action strategy receives the same total addition.
Thus P's shaping can guide learning but cannot, on its own, fix a missing obedience incentive. The same reasoning
covers a report choice only if that choice cannot alter the initial potential. P already invokes prior work on
equilibrium invariance; the peers treated this as a boundary of the connexion, not a new theorem
\cite{devlin2012,lu2014}.

\subsection{A prize and an absolute signal}

The peers' example has two agents. Each chooses effort \(e_i\in\{0,1\}\), and agent \(i\)'s intrinsic payoff is
\(\beta(e_1+e_2)-ke_i\). The number \(\beta\) is the benefit to each agent from one unit of team effort; \(k\) is the
cost of its own effort. The model assumes \(k>\beta\), so effort needs an added incentive. It assumes that agent \(i\)
wins the comparison with probability \(\sigma(\eta(e_i-e_j))\), where \(j\) is the other agent, \(\eta>0\) controls
how informative the comparison is, and \(\sigma(x)=1/(1+\exp(-x))\). Write \(d=\sigma(\eta)-1/2\), the increase in win
probability when one agent alone raises effort.

Now let \(Y_i\) indicate a comparison win and \(Z\) indicate task success. Suppose success has probability
\((e_1+e_2)/2\), and pay agent \(i\) the episode reward \(wY_i+bZ\), with non-negative prize weights \(w\) and \(b\).
Against either action by the other agent, raising effort changes expected payoff by \(\beta-k+wd+b/2\). High effort is
therefore strictly preferable in this model if
\[
wd+b/2>k-\beta.
\]
This is a new episode prize, not P's potential-shaped reward. The calculation appears in the peer check
\texttt{rank-monitor-identification.md}; it requires the stated comparison and success laws even after an agent
changes action or view.

The comparison law is the same at effort profiles \((1,1)\) and \((0,0)\). If both agents switch from high effort to
low effort, each saves \(k\), loses \(2\beta\) in intrinsic benefit, and loses \(b\) in expected success payment. The
joint switch gives each \(k-2\beta-b\). Hence \(b\geq k-2\beta\) blocks this particular joint deviation, while a
comparison-only prize cannot block it when \(k>2\beta\). This coalition check is stronger than Q's unilateral
obedience requirement. P's actual sparse success signal has not been shown to follow the example's assumed probability
law. The general concern about collusion under relative pay predates this thread
\cite{gibbonsmurphy,relative2023,biasedtournaments}.

\subsection{One payment for several deviations}

Ada's \texttt{monitor-obedience.md} check gives an exact test for a fixed target and one agent. Let \(P_0\) be the
probability law of a finite monitor on the target path. Let \(P_j\) be its law after listed unilateral deviation
\(j\), and let \(g_j>0\) be that deviation's expected gain before monitor payments. A single payment based on the
monitor and lying between zero and \(B\) deters all the listed deviations exactly when, for every set of non-negative
weights \(\lambda_j\) summing to one,
\[
B\operatorname{TV}\!\left(P_0,\sum_j\lambda_jP_j\right)
\geq\sum_j\lambda_jg_j.
\]
Here \(\operatorname{TV}\) is total variation distance between probability laws, and \(B\) is the payment cap. The
test considers every mixture of deviations because the same payment must work against each one. Finite minimax
exchanges the choice of payment with the choice of mixture; for a fixed mixture, the largest possible expected payment
difference is \(B\) times its total variation distance from the target law.

The distinction matters. A binary monitor may have a target success probability of \(1/2\) and two deviations with
probabilities zero and one. Each deviation is visible alone, but their equal mixture reproduces the target law. No one
payment deters both if both offer a positive gain. Emmy's \texttt{rank-monitor-identification.md} check also
considered a hypothetical uninformative camera view whose law equals the target law directly. Neither case is an
observation about P's environment. The test does not include Q's full set of false-report and action combinations,
other agents' constraints, or changing coalitions.

\section{Objections and remaining limits}

The peer review narrowed the claim in two ways. First, the telescoping result and the problem of collusion under
relative rewards have prior art, so the effort example is a benchmark rather than a new general result. Second, P's
connected comparison graph alone does not support a finite error bound for its unregularised Bradley--Terry
maximum-likelihood estimate. With two agents and only wins in one direction, the comparison graph is connected, but
increasing the winner's score without bound keeps improving the likelihood. Scores also need a common translation
fixed before an ordinary Euclidean error bound makes sense. P's softmax potential is unchanged by that translation, so
this point need not spoil its shaping rule. The peers' counterexample and the existing maximum-likelihood literature
support the statistical objection \cite{btmle}.

More directly, the prize calculation assumes committed cardinal payments and comparison laws that remain informative
after strategic changes in effort or camera view. P reports errors with uninformative views, but its reported
comparison performance does not measure the laws under deliberate deviations. The finite-monitor test adds a payment
cap \(B\), whereas Q can use richer and unbounded rewards. No empirical result in this thread says whether P's sparse
outcome separates common shirking enough, or whether agents can in fact choose views that erase the comparison signal.

\section{Why the thread stopped}

The verifier paused the thread because its inequalities and monitor test are conditional on hypothetical comparison
and success laws. The peers had not measured those laws under strategic camera choices and shirking in a P-like task.
Without those data, the proposed reward rule does not yet support a paper claim about the pair.

The smallest restart is to choose one two-agent task and list a finite set of shirking and camera-choice deviations.
Sample the comparison and success outputs under the target and each deviation, then solve the bounded-payment
obedience problem for those measured laws. A payment with a positive margin larger than sampling uncertainty would
certify the listed deviations; a profitable mixture with the target monitor law would rule out that monitor and
payment bound.

\bibliographystyle{plainurl}
\bibliography{references}

\end{document}
